\documentclass[11pt]{article}

\usepackage{nmrdoc}
\usepackage{siunitx}
\usepackage{bm}
\usepackage{tikz}

\sisetup{per-mode=symbol}
\DeclareSIUnit{\angstrom}{\text{\AA}}

\newcommand{\code}[1]{\texttt{#1}}
\setlength{\parindent}{0pt}
\setlength{\parskip}{0.42em}
\emergencystretch=2em

\begin{document}
\NmrTitle{Numerical Addendum: McConnell Calculator}

The McConnell calculator is a cutoff sum of closed-form bond kernels. The
numerical work is the bond geometry cached for those kernels, the midpoint
radius query and filters that decide which candidates enter the sum, the
asymmetric per-bond tensor, and the two projection paths used for output. The
listings below are distilled from the source: loops and bookkeeping are
trimmed, but the constants, guards, tensor entries, and output shapes are the
implementation's.

\section{Bond geometry and candidates}

The geometry result builds three arrays indexed by bond: midpoint, length, and
axis. For a bond with endpoint positions \(\bm q_A\) and \(\bm q_B\), the
midpoint is their average, the length is
\(\ell = |\bm q_B-\bm q_A|\), and the axis is the normalized endpoint
difference. If \(\ell\le 10^{-15}\), \code{Bond::Direction} stores the zero
vector. The McConnell kernel consumes that cached vector as it stands; it does
not normalize the bond axis again.

\begin{codebox}
// Distilled from Bond.h and GeometryResult.cpp.
Vec3 edge = positions[bond.atom_index_b] - positions[bond.atom_index_a];
double len = edge.norm();

conf.bond_lengths[bi]    = len;
conf.bond_directions[bi] = (len > 1e-15) ? Vec3(edge / len) : Vec3::Zero();
conf.bond_midpoints[bi]  =
    0.5 * (positions[bond.atom_index_a] + positions[bond.atom_index_b]);
\end{codebox}

For each atom, candidate bonds come from the spatial index over bond
midpoints. The query radius is the runtime TOML key
\code{mcconnell\_bond\_anisotropy\_cutoff}, \(\SI{10.0}{\angstrom}\) in
\code{data/calculator\_params.toml}. The header constant
\code{MCCONNELL\_CUTOFF\_A} is vestigial; the compute path reads the
configuration key shown here. The radius query only chooses candidates. A
returned bond still passes the kernel guard and all three filters before it
changes any accumulator. The cutoff is not a switching function: a bond outside
the query is absent from the per-atom sum, and a returned bond keeps the full
\(r^{-3}\) kernel if the later filters accept it. The per-atom path does not
apply a second cutoff check after this query.

\begin{codebox}
// Distilled from SpatialIndexResult.cpp and McConnellResult.cpp.
auto nearby_bonds =
    spatial.BondsWithinRadius(atom_pos,
        cfg("mcconnell_bond_anisotropy_cutoff"));  // 10.0 A

filters.Add(std::make_unique<MinDistanceFilter>());
filters.Add(std::make_unique<SelfSourceFilter>());
filters.Add(std::make_unique<DipolarNearFieldFilter>());
\end{codebox}

\section{One bond kernel}

For an atom at \(\bm r\), \code{disp} is
\(\bm x=\bm r-\bm m\), the displacement from the bond midpoint to the atom.
Its norm is \(r\), and \code{d\_hat} is
\(\hat{\bm d}=\bm x/r\). The cached \code{bond\_direction} is
\(\bm b\): a unit endpoint axis when \(\ell>10^{-15}\), and the zero vector
when \(\ell\le 10^{-15}\). The code forms one dot product,
\[
  c = \hat{\bm d}\cdot\bm b .
\]
For nondegenerate bonds this is \(\cos\theta\).
The helper returns the zero-initialized result before the division by \(r\)
when \(r<\SI{0.1}{\angstrom}\). There is no clamp to
\(\SI{0.1}{\angstrom}\): below the guard the result's tensor, scalar, distance,
and direction stay at zero. At exactly \(\SI{0.1}{\angstrom}\) the formula is
evaluated with that distance, so the \(r^{-3}\) scale is
\(\SI{1000}{\angstrom^{-3}}\).

\begin{codebox}
// Distilled from ComputeBondKernel.
Vec3 disp = atom_pos - bond_midpoint;      // midpoint -> atom
double r = disp.norm();
if (r < cfg("singularity_guard_distance")) return result;  // zero

result.distance = r;
double r3 = r * r * r;
Vec3 d_hat = disp / r;
double c = d_hat.dot(bond_direction);

result.f = (3.0 * c * c - 1.0) / r3;

for (int a = 0; a < 3; ++a)
    for (int b = 0; b < 3; ++b)
        result.K(a,b) =
            (3.0*d_hat(a)*d_hat(b) - (a == b ? 1.0 : 0.0)) / r3;

for (int a = 0; a < 3; ++a)
    for (int b = 0; b < 3; ++b)
        result.M_over_r3(a,b) =
            ( 9.0*c*d_hat(a)*bond_direction(b)
             -3.0*bond_direction(a)*bond_direction(b)
             -(3.0*d_hat(a)*d_hat(b) - (a == b ? 1.0 : 0.0)) ) / r3;
\end{codebox}

The stored full tensor is
\[
  \mathcal M_{ab} =
  \frac{
    9c\,\hat d_a b_b
    -3 b_a b_b
    -(3\hat d_a\hat d_b-\delta_{ab})}
  {r^3}.
\]
The separately stored dipolar kernel is
\[
  K_{ab}=\frac{3\hat d_a\hat d_b-\delta_{ab}}{r^3}.
\]
Thus the third term of \(\mathcal M\) is exactly \(-K\). For a
nondegenerate bond, the second term is also a symmetric dipolar-shaped dyadic,
but it includes its trace:
\[
  \frac{-3\,\hat{\bm b}\otimes\hat{\bm b}}{r^3}
  =
  -\frac{3\,\hat{\bm b}\otimes\hat{\bm b}-\bm I}{r^3}
  -\frac{\bm I}{r^3}.
\]
After the scalar trace piece is separated, its \(T_2\) part is the negative
bond-axis dipolar kernel. The stored \(K\) is the midpoint-to-atom version of
the same symmetric-traceless shape.

The three terms map to the tensor split in different ways. For a
nondegenerate bond, term 1,
\(9c\,\hat{\bm d}\otimes\hat{\bm b}/r^3\), is generally asymmetric. Its trace
is \(9c^2/r^3\), so it contributes \(3c^2/r^3\) to \(T_0\); its antisymmetric
part is the only source of \(T_1\); and its symmetric traceless remainder
contributes to \(T_2\). Term 2 has trace \(-3/r^3\), so it contributes
\(-1/r^3\) to \(T_0\) and a symmetric-traceless \(T_2\) part, but no \(T_1\).
Term 3 is symmetric traceless, so it contributes only to \(T_2\).
After the distance guard has passed in the degenerate cached-zero-axis case,
term 1 and term 2 vanish, and the code evaluates the scalar formula below with
\(c=0\).

For nondegenerate bonds, those trace pieces explain the scalar. The code
computes
\[
  f=\frac{3c^2-1}{r^3}
\]
directly, and algebraically
\[
  f = \hat b_a K_{ab}\hat b_b = \frac13\operatorname{tr}\mathcal M .
\]
For those bonds, \(f\) is both the double contraction of \(K\) with the bond
axis and the single-bond \(T_0\) value of the full McConnell tensor. It is not
the trace of \(K\), which is zero. When the degenerate cached-zero-axis case is
evaluated, \(f=-1/r^3\) while \(\operatorname{tr}\mathcal M=0\).

For nondegenerate bonds, the same \(c\) threads through three places: it is
squared in \(f\), it is the linear weight on the asymmetric first term, and it
makes the kernel invariant to reversing the bond endpoints. If
\(\hat{\bm b}\) is replaced by \(-\hat{\bm b}\), then \(c\) also changes sign.
The product \(c\,\hat{\bm b}\), the dyadic
\(\hat{\bm b}\otimes\hat{\bm b}\), the stored \(K\), and \(f\) are unchanged.

\begin{figure}[ht]
\centering
\begin{tikzpicture}[scale=1.0,line width=0.45pt,>=stealth]
  \coordinate (A) at (0,0);
  \coordinate (B) at (4.0,0);
  \coordinate (M) at (2.0,0);
  \coordinate (R) at (3.25,2.15);

  \draw (A) -- (B);
  \fill (A) circle (1.5pt) node[below] {\scriptsize \(A\)};
  \fill (B) circle (1.5pt) node[below] {\scriptsize \(B\)};
  \fill (M) circle (1.5pt) node[below] {\scriptsize \(\bm m\)};
  \fill (R) circle (1.5pt) node[right] {\scriptsize atom \(\bm r\)};

  \draw[->] (M) -- (3.05,0)
      node[midway,below] {\scriptsize \(\hat{\bm b}\)};
  \draw[->] (M) -- (R)
      node[midway,left] {\scriptsize \(\hat{\bm d}\)};
  \draw (2.42,0) arc[start angle=0,end angle=59.8,radius=0.42];
  \node at (2.62,0.28) {\scriptsize \(\theta\)};
\end{tikzpicture}
\caption{The bond midpoint supplies the expansion centre. The kernel uses the
midpoint-to-atom direction \(\hat{\bm d}\) and, for nondegenerate bonds, the
cached bond axis \(\hat{\bm b}\); their dot product is \(c=\cos\theta\).}
\end{figure}

\section{Guards and filters}

After the helper returns, the per-atom path builds a
\code{KernelEvaluationContext}. Its \code{distance} field is the helper's
midpoint-to-atom distance, \code{source\_extent} is the cached bond length
\(\ell\), and the two source atom fields are the bond endpoints. The filter set
is ordered as minimum distance, self source, then near field. If any filter
rejects, the candidate leaves the computation at that point: no scalar sum,
tensor sum, bond-neighbour record, nearest-bond record, or category total is
updated.

\begin{codebox}
KernelEvaluationContext ctx;
ctx.distance      = kernel.distance;
ctx.source_extent = conf.bond_lengths[bi];
ctx.atom_index    = ai;
ctx.source_atom_a = bond.atom_index_a;
ctx.source_atom_b = bond.atom_index_b;

if (!filters.AcceptAll(ctx)) continue;
\end{codebox}

\code{MinDistanceFilter} requires
\[
  r \ge \code{singularity\_guard\_distance} = \SI{0.1}{\angstrom}.
\]
For \(r<\SI{0.1}{\angstrom}\), the helper has already returned a zero result
with \code{distance == 0}; this filter keeps that zero-valued placeholder from
being stored as a real neighbour. At \(r=\SI{0.1}{\angstrom}\), the helper has
computed the kernel and this filter accepts it, subject to the later filters.

\code{SelfSourceFilter} is topological. If the target atom is either endpoint
of the bond, the bond is omitted even when the midpoint distance passes the
numeric tests. The filter changes only membership in the sum; it does not alter
the kernel value for other atoms.

\code{DipolarNearFieldFilter} requires
\[
  r >
  \code{near\_field\_exclusion\_ratio}\,\ell
  = 0.5\,\ell .
\]
At or inside half the bond length from the midpoint, the point-source form is
not added to the per-atom sums. Equality is rejected on this path. As with the
other filters, the skipped contribution is not tapered, clamped, or reweighted
into another term.

\code{SampleKernelAt} evaluates the same bond formula at a free point in
space. That path loops over all bonds and has no self-source test because the
point has no atom index. It keeps the same singularity guard, skips
\(r>\SI{10.0}{\angstrom}\), and skips \(r<0.5\ell\). Therefore equality is kept
at the sampling cutoff and at the near-field boundary on the free-point path,
while the per-atom filter path rejects equality at the near-field boundary.

The nearest peptide C=O direction stored on an atom is not the C=O bond axis.
It is \(\hat{\bm d}\), the midpoint-to-atom direction from the nearest accepted
C=O kernel. The code renormalizes that vector before storage; if its norm is
not greater than \code{near\_zero\_vector\_norm\_threshold}, \(10^{-10}\), the
stored direction is the zero vector.

\begin{codebox}
if (best_co_dist < NO_DATA_SENTINEL) {
    double dir_norm = best_co_direction.norm();
    ca.dir_nearest_CO =
        (dir_norm > cfg("near_zero_vector_norm_threshold"))
            ? Vec3(best_co_direction / dir_norm)
            : Vec3::Zero();
    ca.T2_CO_nearest = SphericalTensor::Decompose(best_co_kernel.K);
}
\end{codebox}

\section{Accumulation and projection}

Every accepted bond first contributes its raw asymmetric
\code{M\_over\_r3} to \code{M\_total}. The same raw per-bond tensor may also
enter one broad category total. Backbone C=O, backbone C--N, and other
backbone bonds feed \code{M\_backbone\_total}; side-chain C=O and other
side-chain bonds feed \code{M\_sidechain\_total}; aromatic bonds feed
\code{M\_aromatic\_total}. Disulfide, unknown, or otherwise unhandled
categories still enter \code{M\_total} after the filters, but they do not enter
these broad category totals. The scalar sums are narrower: only peptide C=O,
peptide C--N, side-chain C=O, and aromatic bonds add their \(f\) values.

\begin{codebox}
M_total += kernel.M_over_r3;

case BondCategory::PeptideCO:
    co_sum += kernel.f;
    M_backbone_total += kernel.M_over_r3;
    update_nearest_CO_if_closer();
    break;

case BondCategory::PeptideCN:
    cn_sum += kernel.f;
    M_backbone_total += kernel.M_over_r3;
    update_nearest_CN_if_closer();
    break;

case BondCategory::Aromatic:
    aromatic_sum += kernel.f;
    M_aromatic_total += kernel.M_over_r3;
    break;
\end{codebox}

From this shared accepted-bond input the code takes two different projection
paths. The category tensors first discard the asymmetric part by forming
\(\frac12(M+M^{\mathsf T})\), then remove the trace, then decompose. The
symmetrization discards the \(T_1\) produced by term 1. The trace subtraction
sets \(T_0\) to zero before \code{Decompose} sees the matrix. The packed
category output writes only the five \(T_2\) numbers.

\begin{codebox}
Mat3 S = 0.5 * (M_backbone_total + M_backbone_total.transpose());
S -= (S.trace() / 3.0) * Mat3::Identity();
ca.T2_backbone_total = SphericalTensor::Decompose(S);

S = 0.5 * (M_sidechain_total + M_sidechain_total.transpose());
S -= (S.trace() / 3.0) * Mat3::Identity();
ca.T2_sidechain_total = SphericalTensor::Decompose(S);

S = 0.5 * (M_aromatic_total + M_aromatic_total.transpose());
S -= (S.trace() / 3.0) * Mat3::Identity();
ca.T2_aromatic_total = SphericalTensor::Decompose(S);
\end{codebox}

\code{mc\_shielding\_contribution} takes the other path on the same atom's
accepted raw tensors. It decomposes \code{M\_total} directly, without
symmetrizing and without subtracting the trace first. Its \(T_0\) is the trace
average of the full accepted-bond sum, its \(T_1\) is the antisymmetric part
left by the first dyadic term, and its \(T_2\) is the symmetric traceless part
of that same raw sum. The full nine-column output therefore reports the whole
asymmetric McConnell sum, while the category tensor output reports
symmetric-traceless projections. The nearest C=O and C--N blocks in the
category array are different again: they decompose the stored \(K\) of the
nearest accepted bond of each type, and then write only its \(T_2\) entries.

\begin{codebox}
ca.mc_shielding_contribution =
    SphericalTensor::Decompose(M_total);     // raw M, T0 + T1 + T2

const SphericalTensor* cats[5] = {
    &ca.T2_backbone_total, &ca.T2_sidechain_total,
    &ca.T2_aromatic_total, &ca.T2_CO_nearest, &ca.T2_CN_nearest
};
\end{codebox}

The emitted arrays have fixed second dimensions. If no accepted nearest C=O
or C--N bond is found, the corresponding distance remains
\code{NO\_DATA\_SENTINEL}, \(\SI{99.0}{\angstrom}\).

\begin{codebox}
mc_shielding.npy    // N x 9:  [T0, T1[0..2], T2[0..4]]
mc_category_T2.npy  // N x 25: backbone, sidechain, aromatic, CO, CN
mc_scalars.npy      // N x 6:  four scalar sums, nearest CO/CN distances
\end{codebox}

\section{Shared tensor split}

\code{SphericalTensor::Decompose} is closed-form arithmetic on a \(3\times3\)
matrix \(\bm X\). It does not diagonalize \(\bm X\). The trace gives
\[
  T_0 = \frac{1}{3}(X_{xx}+X_{yy}+X_{zz}).
\]
The antisymmetric part gives the three-vector
\[
  \bm T_1 = \frac12
  \begin{pmatrix}
    X_{yz}-X_{zy}\\
    X_{zx}-X_{xz}\\
    X_{xy}-X_{yx}
  \end{pmatrix}.
\]
The symmetric traceless matrix \(\bm S\) is built as
\[
  S_{ab}=\frac12(X_{ab}+X_{ba})-T_0\delta_{ab}.
\]
The stored five-vector is
\[
  \bm T_2 =
  \begin{pmatrix}
    \sqrt2\,S_{xy}\\
    \sqrt2\,S_{yz}\\
    \sqrt{3/2}\,S_{zz}\\
    \sqrt2\,S_{xz}\\
    (S_{xx}-S_{yy})/\sqrt2
  \end{pmatrix}.
\]

\begin{codebox}
st.T0 = sigma.trace() / 3.0;

st.T1[0] = 0.5 * (sigma(1,2) - sigma(2,1));
st.T1[1] = 0.5 * (sigma(2,0) - sigma(0,2));
st.T1[2] = 0.5 * (sigma(0,1) - sigma(1,0));

double Sxx = sigma(0,0) - st.T0;
double Syy = sigma(1,1) - st.T0;
double Szz = sigma(2,2) - st.T0;
double Sxy = 0.5 * (sigma(0,1) + sigma(1,0));
double Sxz = 0.5 * (sigma(0,2) + sigma(2,0));
double Syz = 0.5 * (sigma(1,2) + sigma(2,1));

st.T2[0] = std::sqrt(2.0) * Sxy;
st.T2[1] = std::sqrt(2.0) * Syz;
st.T2[2] = std::sqrt(3.0 / 2.0) * Szz;
st.T2[3] = std::sqrt(2.0) * Sxz;
st.T2[4] = (Sxx - Syy) / std::sqrt(2.0);
\end{codebox}

The factors on the five \(T_2\) entries make the packing isometric. The
\(\sqrt2\) on an off-diagonal entry accounts for the Frobenius double count:
\(S_{xy}\) appears twice in the matrix, once as \(S_{xy}\) and once as
\(S_{yx}\), so it contributes \(2S_{xy}^2\) to
\(\sum_{ab}S_{ab}^2\). The packed entry \(\sqrt2\,S_{xy}\) squares to that
same number. The diagonal factors account similarly for the two free diagonal
degrees of freedom after \(S_{xx}+S_{yy}+S_{zz}=0\). A dot product of two stored
\(T_2\) five-vectors is therefore the Frobenius inner product of their
symmetric traceless matrices.

McConnell uses all three channels in \code{mc\_shielding\_contribution}. The
dipolar kernel \(K\), the nearest C=O and C--N tensors, and the broad category
tensor outputs populate the \(T_2\) channel by construction; their \(T_0\) and
\(T_1\) entries are not written to \code{mc\_category\_T2.npy}.

\section{Glossary}

\textbf{Dyadic product.} A grid whose rows are labeled by one vector and whose
columns are labeled by another, with each cell holding the product of its two
labels. In the nondegenerate McConnell kernel,
\(\hat{\bm d}\otimes\hat{\bm b}\) has entry \(\hat d_a\hat b_b\), and it is
asymmetric unless the two directions are parallel or antiparallel. Formally,
for column vectors \(\bm u,\bm v\),
\((\bm u\otimes\bm v)_{ab}=u_a v_b\).

\textbf{Dipolar kernel.} The field shape of a small bar magnet --- strongest
and signed along its axis, weaker and reversed to the side, fading as the cube
of distance (an axially symmetric, trace-zero tensor). Along its axis it has
eigenvalue \(+2/r^3\), while either broadside direction has
eigenvalue \(-1/r^3\); the trace is zero because those three directions
balance. Formally, for an axis \(\hat{\bm u}\), it is the symmetric traceless
tensor \((3\hat u_a\hat u_b-\delta_{ab})/r^3\).

\textbf{Symmetric-traceless projection.} Drop a matrix's rotational swirl and
its uniform scaling, and what is left is pure shape (remove the antisymmetric
part and the trace average; the result is symmetric and has zero trace). The
broad McConnell category totals use it by forming \(0.5(M+M^{\mathsf T})\), then
subtracting one third of that trace from each diagonal entry before
decomposition. Formally,
\[
  \operatorname{STF}(M)=\frac12(M+M^{\mathsf T})
  -\frac13\operatorname{tr}\!\left(\frac12(M+M^{\mathsf T})\right)\bm I .
\]

\textbf{Near-field exclusion.} A bubble drawn around the source, sized to the
real object the point kernel stands in for (an exclusion region scaled to that
object). For a
bond of length \(\ell\), the per-atom McConnell path omits the candidate unless
the midpoint distance is greater than \(0.5\ell\); equality is also omitted on
that path. Formally, it is the filter \(r>\rho_{\mathrm{nf}}\ell\) with
\(\rho_{\mathrm{nf}}=\code{near\_field\_exclusion\_ratio}=0.5\).

\textbf{Spherical tensor split.} A \(3\times3\) tensor holds three kinds of
content that rotate without mixing --- an orientation-free size (the scalar
trace, \(T_0\)), an axial twist from its antisymmetric part (the three \(T_1\)
components), and a five-number traceless shape for the rest of the
direction-dependence (the five \(T_2\) components). Here the full McConnell sum
can fill all nine stored numbers, while \(K\) and the category projections
contribute only the five \(T_2\) numbers that are written in the category array.
Formally, it decomposes the nine Cartesian components into dimensions
\(1+3+5\), with the \(T_2\) basis normalized to preserve the Frobenius norm.

\end{document}
